% ===============================================================
% WEEKLY REPORT — copy this file to week02.tex, week03.tex, ...
% In Overleaf: use the file dropdown above the editor (or
% Menu > Settings > Main document) to pick which root to compile.
% ===============================================================
\documentclass[11pt,a4paper]{article}

\input{setup/project}
\input{setup/preamble}
\input{setup/weekly}

% No paragraph indentation; separate paragraphs with vertical space instead.
\setlength{\parindent}{0pt}
\setlength{\parskip}{0.6em}


\begin{document}

\weeklytitle{3}

\section{Week 3: 14--18 September 2026}

\sectionprogress

\textbf{Coronary tree graph built from the centerlines.} From the raw centerlines, points and
polyline connections per vessel side, \texttt{topology/graph.py} builds a tree that knows the
ostium (root), which way is downstream, and which segments chain into one continuous named vessel
(LAD, LCx, etc.). This matters for every feature built on top: a bifurcation angle only means
something once each branch's direction points away from the junction, and a vessel's length or
tortuosity only accumulates correctly walking ostium to tip. I checked the tree against the
delivered flags marking "branch point" / "endpoint."

\textbf{First features: angle, volume, tortuosity.} On top of that tree I started computing
bifurcation angle, angle together with local branch volume, and tortuosity, split into its several
distinct parts: the straight-line-distance ratio ($L/D$), local bending (curvature), out-of-plane
corkscrewing (torsion, non-planarity), and how often the vessel oscillates (inflections, SOAM).

\textbf{Tortuosity's sharper measures may not be trustworthy.} Curvature needs a second derivative
of the centerline and torsion a third; derivatives amplify noise, so these two measures may be far
less reliable than the others.

\textbf{CAS-Net pipeline reproduced end to end.} I went through the whole pipeline and wrote up
each stage, trained from scratch for 5 epochs, then ran the full inference and evaluation off the
delivered pretrained weights and reproduced the expected results.

\include{content/methods/topology/volume_disease}
\include{content/methods/topology/tortuosity}
\include{content/methods/pipeline/data_preparation}
\include{content/methods/pipeline/model_architecture}
\include{content/methods/pipeline/training}
\include{content/methods/pipeline/inference}
\include{content/methods/pipeline/evaluation_metrics}



\end{document}